\input{preamble}

% OK, start here.
%
\begin{document}

\title{Exemples de champs}


\maketitle

\phantomsection
\label{section-phantom}

\tableofcontents

\section{Introduction}
\label{section-introduction}

\noindent
Nous étudions ici des exemples de champs en géométrie algébrique.
Certains sont des champs algébriques, d'autres ne le sont pas.
Nous déterminerons lesquels sont des champs algébriques dans un chapitre ultérieur.
Dans le présent chapitre, nous nous préoccupons donc principalement des conditions
de descente. Voir par exemple \cite{Vis2}.

\medskip\noindent
Certaines notations, conventions et certains termes de ce chapitre sont peu naturels
et pourront sembler inversés au lecteur plus expérimenté. Cela est intentionnel.
On trouvera une explication dans Quot, section \ref{quot-section-conventions}.




\section{Notations}
\label{section-notation}

\noindent
Dans ce chapitre, nous fixons un grand site fppf convenable $\Sch_{fppf}$
comme dans Topologies, définition \ref{topologies-definition-big-fppf-site}.
Ainsi, sauf indication explicite contraire, tous les schémas seront des objets
de $\Sch_{fppf}$.
Nous travaillerons toujours relativement à une base $S$ contenue dans $\Sch_{fppf}$.
Nous travaillerons alors avec le grand site fppf $(\Sch/S)_{fppf}$,
voir Topologies, définition \ref{topologies-definition-big-small-fppf}.
On retrouve le cas absolu en prenant
$S = \Spec(\mathbf{Z})$.




\section{Exemples de champs}
\label{section-examples-stacks}

\noindent
Nous donnons d'abord quelques exemples importants de champs sur
$(\Sch/S)_{fppf}$.




\section{Faisceaux quasi-cohérents}
\label{section-stack-of-quasi-coherent-sheaves}

\noindent
Nous définissons une catégorie $\QCohstack$ comme suit :
\begin{enumerate}
\item Un objet de $\QCohstack$ est un couple $(X, \mathcal{F})$,
où $X/S$ est un objet de $(\Sch/S)_{fppf}$ et $\mathcal{F}$
est un $\mathcal{O}_X$-module quasi-cohérent ;
\item un morphisme $(f, \varphi) : (Y, \mathcal{G}) \to (X, \mathcal{F})$
est un couple formé d'un morphisme $f : Y \to X$ de schémas sur $S$
et d'un $f$-morphisme (voir
Faisceaux, section \ref{sheaves-section-ringed-spaces-functoriality-modules})
$\varphi : \mathcal{F} \to \mathcal{G}$ ;
\item le composé des morphismes
$$
(Z, \mathcal{H}) \xrightarrow{(g, \psi)}
(Y, \mathcal{G}) \xrightarrow{(f, \phi)} (X, \mathcal{F})
$$
est $(f \circ g, \psi \circ \phi)$, où $\psi \circ \phi$ est
le composé des $f$-morphismes.
\end{enumerate}
Ainsi, $\QCohstack$ est une catégorie et
$$
p : \QCohstack \to (\Sch/S)_{fppf},
\quad
(X, \mathcal{F}) \mapsto X
$$
est un foncteur. Remarquons que la catégorie fibre de $\QCohstack$ au-dessus
d'un schéma $X$ est l'opposée de la catégorie $\QCoh(\mathcal{O}_X)$
des $\mathcal{O}_X$-modules quasi-cohérents.
Pour un usage ultérieur, remarquons qu'étant donnés
$(X, \mathcal{F}), (Y, \mathcal{G}) \in \Ob(\QCohstack)$,
nous avons
\begin{equation}
\label{equation-morphisms-qcoh}
\Mor_{\QCohstack}((Y, \mathcal{G}), (X, \mathcal{F}))
=
\coprod\nolimits_{f \in \Mor_S(Y, X)}
\Mor_{\QCoh(\mathcal{O}_Y)}(f^*\mathcal{F}, \mathcal{G}).
\end{equation}
Voir la discussion des $f$-morphismes de modules dans
Faisceaux, section \ref{sheaves-section-ringed-spaces-functoriality-modules}.

\medskip\noindent
La catégorie $\QCohstack$ n'est pas un champ sur $(\Sch/S)_{fppf}$,
car sa collection d'objets est une classe propre. Nous verrons néanmoins
qu'elle satisfait tous les axiomes d'un champ. Nous contournerons
la difficulté ensembliste dans la
section \ref{section-stack-of-finitely-generated-quasi-coherent-sheaves}.

\begin{lemma}
\label{lemma-quasi-coherent-strongly-cartesian}
Un morphisme $(f, \varphi) : (Y, \mathcal{G}) \to (X, \mathcal{F})$
de $\QCohstack$ est fortement cartésien si et seulement si le
morphisme $\varphi$ induit un isomorphisme $f^*\mathcal{F} \to \mathcal{G}$.
\end{lemma}

\begin{proof}
Soit $(X, \mathcal{F}) \in \Ob(\QCohstack)$.
Soit $f : Y \to X$ un morphisme de $(\Sch/S)_{fppf}$.
Il existe un $f$-morphisme canonique $c : \mathcal{F} \to f^*\mathcal{F}$
et nous obtenons donc un morphisme
$(f, c) : (Y, f^*\mathcal{F}) \to (X, \mathcal{F})$.
Nous affirmons que $(f, c)$ est fortement cartésien.
En effet, pour tout objet $(Z, \mathcal{H})$ de $\QCohstack$, nous avons
\begin{align*}
\Mor_{\QCohstack}((Z, \mathcal{H}), (Y, f^*\mathcal{F}))
& =
\coprod\nolimits_{g \in \Mor_S(Z, Y)}
\Mor_{\QCoh(\mathcal{O}_Z)}(g^*f^*\mathcal{F}, \mathcal{H}) \\
& =
\coprod\nolimits_{g \in \Mor_S(Z, Y)}
\Mor_{\QCoh(\mathcal{O}_Z)}((f \circ g)^*\mathcal{F}, \mathcal{H}) \\
& =
\Mor_{\QCohstack}((Z, \mathcal{H}), (X, \mathcal{F}))
\times_{\Mor_S(Z, X)} \Mor_S(Z, Y),
\end{align*}
où nous avons utilisé deux fois l'équation (\ref{equation-morphisms-qcoh}).
Cela prouve que la condition de
Catégories, définition \ref{categories-definition-cartesian-over-C},
est satisfaite pour $(f, c)$, d'où notre assertion. Par
Catégories, lemme \ref{categories-lemma-composition-cartesian},
les isomorphismes sont fortement cartésiens et
les composés de morphismes fortement cartésiens le sont aussi,
ce qui prouve l'implication « si » du lemme. Réciproquement,
étant donnés $(X, \mathcal{F})$ et $f : Y \to X$, s'il existe un
morphisme fortement cartésien relevant $f$ et de but $(X, \mathcal{F})$,
alors il est nécessairement isomorphe à $(f, c)$ (voir la discussion qui suit
Catégories, définition \ref{categories-definition-cartesian-over-C}).
L'implication « seulement si » du lemme en résulte.
\end{proof}

\begin{lemma}
\label{lemma-stack-of-quasi-coherent-sheaves}
Le foncteur $p : \QCohstack \to (\Sch/S)_{fppf}$
satisfait les conditions (1), (2) et (3) de
Champs, définition \ref{stacks-definition-stack}.
\end{lemma}

\begin{proof}
Il résulte clairement du
lemme \ref{lemma-quasi-coherent-strongly-cartesian}
que $\QCohstack$ est une catégorie fibrée sur $(\Sch/S)_{fppf}$.
Pour un recouvrement $\mathcal{U} = \{X_i \to X\}_{i \in I}$ de
$(\Sch/S)_{fppf}$, le foncteur
$$
\QCoh(\mathcal{O}_X) \longrightarrow DD(\mathcal{U})
$$
est pleinement fidèle et essentiellement surjectif, voir
Descente, proposition \ref{descent-proposition-fpqc-descent-quasi-coherent}.
Par conséquent,
Champs, lemme \ref{stacks-lemma-stack-equivalences},
montre que $\QCohstack$ satisfait tous les
axiomes d'un champ.
\end{proof}





\section{Le champ des faisceaux quasi-cohérents de type fini}
\label{section-stack-of-finitely-generated-quasi-coherent-sheaves}

\noindent
On peut en fait obtenir un champ de faisceaux quasi-cohérents
si l'on ne considère que les modules quasi-cohérents de type fini.
Notons
$$
p_{fg} : \QCohstack_{fg} \to (\Sch/S)_{fppf}
$$
la sous-catégorie pleine de $\QCohstack$ sur $(\Sch/S)_{fppf}$
formée des couples $(T, \mathcal{F})$ tels que $\mathcal{F}$
soit un $\mathcal{O}_T$-module quasi-cohérent de type fini.

\begin{lemma}
\label{lemma-stack-of-finite-type-quasi-coherent-sheaves}
Le foncteur $p_{fg} : \QCohstack_{fg} \to (\Sch/S)_{fppf}$
satisfait les conditions (1), (2) et (3) de
Champs, définition \ref{stacks-definition-stack}.
\end{lemma}

\begin{proof}
Nous allons pour cela vérifier les hypothèses (1), (2) et (3) de
Champs, lemme \ref{stacks-lemma-substack}.
D'après le
lemme \ref{lemma-quasi-coherent-strongly-cartesian},
un morphisme $(Y, \mathcal{G}) \to (X, \mathcal{F})$ est
fortement cartésien si et seulement s'il induit un isomorphisme
$f^*\mathcal{F} \to \mathcal{G}$. D'après
Modules, lemme \ref{modules-lemma-pullback-finite-type},
l'image inverse d'un $\mathcal{O}_X$-module de type fini est de type
fini. L'hypothèse (1) de
Champs, lemme \ref{stacks-lemma-substack},
est donc satisfaite. L'hypothèse (2) est immédiate.
Enfin, pour prouver l'hypothèse (3), il faut montrer ceci :
si $\mathcal{F}$ est un $\mathcal{O}_X$-module quasi-cohérent
et si $\{f_i : X_i \to X\}$ est un recouvrement fppf tel que chaque
$f_i^*\mathcal{F}$ soit de type fini, alors $\mathcal{F}$ est de
type fini. En restreignant $\mathcal{F}$ à
un ouvert affine de $X$, on se ramène à l'énoncé algébrique suivant :
soient $R \to S$ un homomorphisme d'anneaux de présentation finie et fidèlement plat,
et $M$ un $R$-module. Si $M \otimes_R S$ est un
$S$-module de type fini, alors $M$ est un $R$-module de type fini.
Une forme plus forte de ce fait algébrique se trouve dans
Algèbre, lemme \ref{algebra-lemma-descend-properties-modules}.
\end{proof}

\begin{lemma}
\label{lemma-finite-type}
Soit $(X, \mathcal{O}_X)$ un espace annelé.
\begin{enumerate}
\item La catégorie des $\mathcal{O}_X$-modules de type fini possède un
ensemble de classes d'isomorphisme.
\item La catégorie des $\mathcal{O}_X$-modules quasi-cohérents
de type fini possède un ensemble de classes d'isomorphisme.
\end{enumerate}
\end{lemma}

\begin{proof}
La partie (2) résulte de la partie (1), car la catégorie de (2) est une sous-catégorie
pleine de celle de (1). Considérons un recouvrement ouvert quelconque
$\mathcal{U} : X = \bigcup_{i \in I} U_i$. Notons $j_i : U_i \to X$
les morphismes d'inclusion. Considérons une application $r : I \to \mathbf{N}$.
Si $\mathcal{F}$ est un $\mathcal{O}_X$-module dont la restriction à
$U_i$ est engendrée par au plus $r(i)$ sections de $\mathcal{F}(U_i)$,
alors $\mathcal{F}$ est un quotient du faisceau
$$
\mathcal{H}_{\mathcal{U}, r} =
\bigoplus\nolimits_{i \in I} j_{i, !}\mathcal{O}_{U_i}^{\oplus r(i)}.
$$
Par définition, si $\mathcal{F}$ est de type fini, il existe
un recouvrement ouvert $\mathcal{U}$ dont l'ensemble d'indices est $I = X$
et pour lequel cette condition est vérifiée. Il suffit donc de montrer qu'il existe
un ensemble de choix possibles pour $\mathcal{U}$ (c'est évident),
un ensemble de choix possibles pour $r : I \to \mathbf{N}$ (c'est évident) et,
pour chaque $\mathcal{U}$ et chaque $r$, un ensemble de modules quotients possibles
de $\mathcal{H}_{\mathcal{U}, r}$. Autrement dit, il suffit de montrer
que, pour un $\mathcal{O}_X$-module $\mathcal{H}$ donné, les quotients ont
au plus un ensemble de classes d'isomorphisme.
Cette dernière assertion devient évidente
si l'on considère les noyaux d'un morphisme quotient
$\mathcal{H} \to \mathcal{F}$
comme paramétrés par une partie de l'ensemble des parties de
$\prod_{U \subset X\text{ ouvert}} \mathcal{H}(U)$.
\end{proof}

\begin{lemma}
\label{lemma-stack-fg-quasi-coherent}
Il existe une sous-catégorie
$\QCohstack_{fg, small} \subset \QCohstack_{fg}$
ayant les propriétés suivantes :
\begin{enumerate}
\item le foncteur d'inclusion
$\QCohstack_{fg, small} \to \QCohstack_{fg}$ est
pleinement fidèle et essentiellement surjectif ;
\item le foncteur
$p_{fg, small} : \QCohstack_{fg, small} \to (\Sch/S)_{fppf}$
fait de $\QCohstack_{fg, small}$ un champ sur $(\Sch/S)_{fppf}$.
\end{enumerate}
\end{lemma}

\begin{proof}
Nous avons vu dans les
lemmes \ref{lemma-stack-of-finite-type-quasi-coherent-sheaves} et
\ref{lemma-finite-type}
que $p_{fg} : \QCohstack_{fg} \to (\Sch/S)_{fppf}$
satisfait les conditions (1), (2) et (3) de
Champs, définition \ref{stacks-definition-stack},
ainsi que la condition supplémentaire (4) de
Champs, remarque \ref{stacks-remark-stack-make-small}.
La discussion de cette remarque fournit donc $\QCohstack_{fg, small}$.
\end{proof}

\noindent
Nous effectuerons souvent le remplacement
$$
\QCohstack_{fg} \leadsto \QCohstack_{fg, small}
$$
sans le signaler davantage et, par abus de notation, nous désignerons
simplement ce remplacement par $\QCohstack_{fg}$.

\begin{remark}
\label{remark-higher-rank}
Toute la discussion de cette section reste valable
si l'on veut considérer les
faisceaux quasi-cohérents qui sont localement engendrés par au plus $\kappa$
sections, pour un cardinal infini $\kappa$, par exemple $\kappa = \aleph_0$.
\end{remark}





\section{Revêtements étales finis}
\label{section-finite-etale}

\noindent
Nous définissons une catégorie $\textit{F\'Et}$ comme suit :
\begin{enumerate}
\item Un objet de $\textit{F\'Et}$ est un morphisme étale fini $Y \to X$
de schémas (d'après nos conventions, il s'agit d'un morphisme étale fini
dans $(\Sch/S)_{fppf}$) ;
\item un morphisme $(b, a) : (Y \to X) \to (Y' \to X')$ de $\textit{F\'Et}$
est un diagramme commutatif
$$
\xymatrix{
Y \ar[d] \ar[r]_b & Y' \ar[d] \\
X \ar[r]_a & X'
}
$$
dans la catégorie des schémas.
\end{enumerate}
Ainsi, $\textit{F\'Et}$ est une catégorie et
$$
p : \textit{F\'Et} \to (\Sch/S)_{fppf},
\quad
(Y \to X) \mapsto X
$$
est un foncteur. La catégorie fibre de $\textit{F\'Et}$ au-dessus
d'un schéma $X$ est précisément la catégorie $\textit{F\'Et}_X$ étudiée dans
Groupes fondamentaux, section \ref{pione-section-finite-etale}.

\begin{lemma}
\label{lemma-finite-etale-stack}
Le foncteur
$$
p : \textit{F\'Et} \longrightarrow (\Sch/S)_{fppf}
$$
définit un champ sur $(\Sch/S)_{fppf}$.
\end{lemma}

\begin{proof}
La descente fppf des morphismes étales finis résulte de
Descente, lemmes \ref{descent-lemma-affine},
\ref{descent-lemma-descending-property-finite} et
\ref{descent-lemma-descending-property-etale}.
Nous omettons les détails.
\end{proof}






\section{Espaces algébriques}
\label{section-stack-of-spaces}

\noindent
Nous définissons une catégorie $\Spacesstack$ comme suit :
\begin{enumerate}
\item Un objet de $\Spacesstack$ est un morphisme $X \to U$
d'espaces algébriques sur $S$, où $U$ est représentable par un objet de
$(\Sch/S)_{fppf}$ ;
\item un morphisme $(f, g) : (X \to U) \to (Y \to V)$
est un diagramme commutatif
$$
\xymatrix{
X \ar[d] \ar[r]_f & Y \ar[d] \\
U \ar[r]^g & V
}
$$
de morphismes d'espaces algébriques sur $S$.
\end{enumerate}
Ainsi, $\Spacesstack$ est une catégorie et
$$
p : \Spacesstack \to (\Sch/S)_{fppf},
\quad
(X \to U) \mapsto U
$$
est un foncteur. La catégorie fibre de $\Spacesstack$ au-dessus
d'un schéma $U$ est précisément la catégorie $\textit{Espaces}/U$ des
espaces algébriques sur $U$ (voir
Topologies sur les espaces, section \ref{spaces-topologies-section-procedure}).
Nous considérerons donc parfois un objet de $\Spacesstack$ comme un
couple $X/U$ formé d'un schéma $U$ et d'un espace algébrique $X$ sur $U$.
Pour un usage ultérieur, remarquons qu'étant donnés
$(X/U), (Y/V) \in \Ob(\Spacesstack)$,
nous avons
\begin{equation}
\label{equation-morphisms-spaces}
\Mor_{\Spacesstack}(X/U, Y/V)
=
\coprod\nolimits_{g \in \Mor_S(U, V)}
\Mor_{\textit{Espaces}/U}(X, U \times_{g, V} Y).
\end{equation}
La catégorie $\Spacesstack$ est presque, mais pas tout à fait, un champ
sur $(\Sch/S)_{fppf}$. La difficulté est de nature ensembliste,
comme nous l'expliquerons ci-dessous.

\begin{lemma}
\label{lemma-spaces-strongly-cartesian}
Un morphisme $(f, g) : X/U \to Y/V$
de $\Spacesstack$ est fortement cartésien si et seulement si le
morphisme $f$ induit un isomorphisme $X \to U \times_{g, V} Y$.
\end{lemma}

\begin{proof}
Soit $Y/V \in \Ob(\Spacesstack)$.
Soit $g : U \to V$ un morphisme de $(\Sch/S)_{fppf}$.
La projection $p : U \times_{g, V} Y \to Y$
donne un morphisme
$(p, g) : U \times_{g, V} Y/U \to Y/V$ de $\Spacesstack$.
Nous affirmons que $(p, g)$ est fortement cartésien.
En effet, pour tout objet $Z/W$ de $\Spacesstack$, nous avons
\begin{align*}
\Mor_{\Spacesstack}(Z/W, U \times_{g, V} Y/U)
& =
\coprod\nolimits_{h \in \Mor_S(W, U)}
\Mor_{\textit{Espaces}/W}(Z, W \times_{h, U} U \times_{g, V} Y) \\
& =
\coprod\nolimits_{h \in \Mor_S(W, U)}
\Mor_{\textit{Espaces}/W}(Z, W \times_{g \circ h, V} Y) \\
& =
\Mor_{\Spacesstack}(Z/W, Y/V)
\times_{\Mor_S(W, V)} \Mor_S(W, U),
\end{align*}
où nous avons utilisé deux fois l'équation (\ref{equation-morphisms-spaces}).
Cela prouve que la condition de
Catégories, définition \ref{categories-definition-cartesian-over-C},
est satisfaite pour $(p, g)$, d'où notre assertion. Par
Catégories, lemme \ref{categories-lemma-composition-cartesian},
les isomorphismes sont fortement cartésiens et
les composés de morphismes fortement cartésiens le sont aussi,
ce qui prouve l'implication « si » du lemme. Réciproquement,
étant donnés $Y/V$ et $g : U \to V$, s'il existe un
morphisme fortement cartésien relevant $g$ et de but $Y/V$,
alors il est nécessairement isomorphe à $(p, g)$ (voir la discussion qui suit
Catégories, définition \ref{categories-definition-cartesian-over-C}).
L'implication « seulement si » du lemme en résulte.
\end{proof}

\begin{lemma}
\label{lemma-pre-stack-of-spaces}
Le foncteur $p : \Spacesstack \to (\Sch/S)_{fppf}$
satisfait les conditions (1) et (2) de
Champs, définition \ref{stacks-definition-stack}.
\end{lemma}

\begin{proof}
Il résulte du
lemme \ref{lemma-spaces-strongly-cartesian}
que $\Spacesstack$ est une catégorie fibrée sur $(\Sch/S)_{fppf}$,
ce qui prouve (1).
Supposons que $\{U_i \to U\}_{i \in I}$ soit un recouvrement de
$(\Sch/S)_{fppf}$. Soient $X, Y$ des espaces algébriques sur
$U$. Enfin, supposons donnés des morphismes $\varphi_i : X_{U_i} \to Y_{U_i}$
de $\textit{Espaces}/U_i$ tels que les restrictions de $\varphi_i$ et $\varphi_j$
soient les mêmes morphismes $X_{U_i \times_U U_j} \to Y_{U_i \times_U U_j}$
d'espaces algébriques sur $U_i \times_U U_j$.
Pour prouver (2), il faut montrer qu'il existe un unique morphisme
$\varphi  : X \to Y$ sur $U$ dont le changement de base à $U_i$ est
égal à $\varphi_i$. Comme un morphisme de $X$ vers $Y$ équivaut à
un morphisme de faisceaux, cela résulte directement de
Sites, lemme \ref{sites-lemma-glue-maps}.
\end{proof}

\begin{remark}
\label{remark-stack-spaces}
Si l'on néglige les difficultés ensemblistes\footnote{La difficulté ne tient pas
à ce que $\Spacesstack$ soit une classe propre, puisque notre définition d'un
espace algébrique sur $S$ n'autorise qu'un ensemble de
classes d'isomorphisme d'espaces algébriques sur $S$. Elle tient plutôt à ce que
des réunions disjointes arbitraires d'espaces algébriques puissent devenir trop grandes
et sortir ainsi de l'« univers partiel » d'ensembles que nous avons choisi.},
$\Spacesstack$ satisfait aussi
la descente des objets et est donc un champ. Il faut en effet montrer qu'étant donnés
\begin{enumerate}
\item un recouvrement fppf $\{U_i \to U\}_{i \in I}$ ;
\item pour chaque $i \in I$, un espace algébrique $X_i/U_i$ ;
\item pour chaque $i, j \in I$, un isomorphisme
$\varphi_{ij} : X_i \times_U U_j \to U_i \times_U X_j$ d'espaces algébriques
sur $U_i \times_U U_j$, satisfaisant la condition de cocycle sur
$U_i \times_U U_j \times_U U_k$,
\end{enumerate}
il existe un espace algébrique $X/U$ et des isomorphismes
$X_{U_i} \cong X_i$ sur $U_i$ qui redonnent les isomorphismes $\varphi_{ij}$.
D'abord, d'après
Sites, lemme \ref{sites-lemma-glue-sheaves},
il existe sur $(\Sch/U)_{fppf}$ un faisceau $X$ qui redonne
les $X_i$ et les $\varphi_{ij}$. Puis, d'après
Amorçage, lemme \ref{bootstrap-lemma-locally-algebraic-space},
$X$ est un espace algébrique (si l'on néglige la condition ensembliste
de ce lemme).
Nous emploierons cet argument dans la section suivante pour montrer que,
si l'on ne considère que les espaces algébriques de type fini, on obtient
un champ.
\end{remark}





\section{Le champ des espaces algébriques de type fini}
\label{section-stack-of-finite-type-spaces}


\noindent
On peut en fait obtenir un champ d'espaces
si l'on ne considère que les espaces de type fini.
Notons
$$
p_{ft} : \Spacesstack_{ft} \to (\Sch/S)_{fppf}
$$
la sous-catégorie pleine de $\Spacesstack$ sur $(\Sch/S)_{fppf}$
formée des couples $X/U$ tels que $X \to U$
soit un morphisme de type fini.

\begin{lemma}
\label{lemma-stack-of-finite-type-spaces}
Le foncteur
$p_{ft} : \Spacesstack_{ft} \to (\Sch/S)_{fppf}$
satisfait les conditions (1), (2) et (3) de
Champs, définition \ref{stacks-definition-stack}.
\end{lemma}

\begin{proof}
Nous allons donner tous les détails, jusqu'à l'excès (ce qui risque
de masquer l'idée de la démonstration).

\medskip\noindent
Nous avons vu dans le
lemme \ref{lemma-spaces-strongly-cartesian}
qu'un morphisme $(f, g) : X/U \to Y/V$ de $\Spacesstack$ est
fortement cartésien si le morphisme induit $f : X \to U \times_V Y$
est un isomorphisme. Si $Y \to V$ est de type fini,
alors $U \times_V Y \to U$ l'est aussi, voir
Morphismes d'espaces,
lemme \ref{spaces-morphisms-lemma-base-change-finite-type}.
Ainsi, si $(f, g) : X/U \to Y/V$ dans $\Spacesstack$ est
fortement cartésien dans $\Spacesstack$ et si $Y/V$ est un objet
de $\Spacesstack_{ft}$, alors $X/U$ est automatiquement aussi un
objet de $\Spacesstack_{ft}$ et, bien entendu, $(f, g)$ est
également fortement cartésien dans $\Spacesstack_{ft}$. On en
conclut que $\Spacesstack_{ft}$ est une catégorie fibrée sur
$(\Sch/S)_{fppf}$. Cela prouve (1).

\medskip\noindent
L'argument précédent montre aussi que le foncteur d'inclusion
$\Spacesstack_{ft} \to \Spacesstack$ transforme les
morphismes fortement cartésiens en morphismes fortement cartésiens.
Autrement dit, $\Spacesstack_{ft} \to \Spacesstack$ est
un $1$-morphisme de catégories fibrées sur $(\Sch/S)_{fppf}$.

\medskip\noindent
Soit $U \in \Ob((\Sch/S)_{fppf})$.
Soient $X, Y$ des espaces algébriques de type fini sur $U$. D'après
Champs, lemme \ref{stacks-lemma-presheaf-mor-map-fibred-categories},
nous obtenons un morphisme de préfaisceaux
$$
\mathit{Mor}_{\Spacesstack_{ft}}(X, Y)
\longrightarrow
\mathit{Mor}_{\Spacesstack}(X, Y)
$$
qui est un isomorphisme, puisque $\Spacesstack_{ft}$ est une sous-catégorie pleine de
$\Spacesstack$. Le membre de gauche est donc un faisceau, car nous avons montré,
au lemme \ref{lemma-pre-stack-of-spaces},
que celui de droite en est un. Cela prouve (2).

\medskip\noindent
Pour prouver la condition (3) de
Champs, définition \ref{stacks-definition-stack},
il faut montrer ceci : étant donnés
\begin{enumerate}
\item un recouvrement $\{U_i \to U\}_{i \in I}$ de $(\Sch/S)_{fppf}$ ;
\item pour chaque $i \in I$, un espace algébrique $X_i$ de type fini sur $U_i$ ;
\item pour chaque $i, j \in I$, un isomorphisme
$\varphi_{ij} : X_i \times_U U_j \to U_i \times_U X_j$ d'espaces algébriques
sur $U_i \times_U U_j$, satisfaisant la condition de cocycle sur
$U_i \times_U U_j \times_U U_k$,
\end{enumerate}
il existe un espace algébrique $X$ de type fini sur $U$ et des isomorphismes
$X_{U_i} \cong X_i$ sur $U_i$ qui redonnent les isomorphismes $\varphi_{ij}$.
Cela résulte de
Amorçage, lemme \ref{bootstrap-lemma-descend-algebraic-space}, partie (2). D'après
Descente sur les espaces, lemme
\ref{spaces-descent-lemma-descending-property-locally-finite-presentation},
le morphisme $X \to U$ est de type fini, ce qui achève la démonstration.
\end{proof}

\begin{lemma}
\label{lemma-stack-ft-spaces}
Il existe une sous-catégorie
$\Spacesstack_{ft, small} \subset \Spacesstack_{ft}$
ayant les propriétés suivantes :
\begin{enumerate}
\item le foncteur d'inclusion
$\Spacesstack_{ft, small} \to \Spacesstack_{ft}$ est
pleinement fidèle et essentiellement surjectif ;
\item le foncteur
$p_{ft, small} : \Spacesstack_{ft, small} \to (\Sch/S)_{fppf}$
fait de $\Spacesstack_{ft, small}$ un champ sur
$(\Sch/S)_{fppf}$.
\end{enumerate}
\end{lemma}

\begin{proof}
Nous avons vu dans le
lemme \ref{lemma-stack-of-finite-type-spaces}
que $p_{ft} : \Spacesstack_{ft} \to (\Sch/S)_{fppf}$
satisfait les conditions (1), (2) et (3) de
Champs, définition \ref{stacks-definition-stack}.
La condition supplémentaire (4) de
Champs, remarque \ref{stacks-remark-stack-make-small},
est satisfaite parce que tout espace algébrique $X$ sur $S$ est de la
forme $U/R$, avec $U, R \in \Ob((\Sch/S)_{fppf})$, voir
Espaces, lemme \ref{spaces-lemma-space-presentation}.
Il n'existe donc qu'un ensemble de classes d'isomorphisme d'objets.
La discussion de cette remarque fournit alors $\Spacesstack_{ft, small}$.
\end{proof}

\noindent
Nous effectuerons souvent le remplacement
$$
\Spacesstack_{ft} \leadsto \Spacesstack_{ft, small}
$$
sans le signaler davantage et, par abus de notation, nous désignerons
simplement ce remplacement par $\Spacesstack_{ft}$.

\begin{remark}
\label{remark-higher-cardinality-spaces}
Toute la discussion de cette section reste valable
si l'on veut considérer les espaces algébriques $X/U$ qui sont
localement de type fini et tels que l'image inverse dans $X$ d'un ouvert affine
de $U$ puisse être recouverte par un nombre dénombrable d'ouverts affines.
Au besoin, on peut aussi introduire la notion de morphisme de
$\kappa$-type (c'est-à-dire imposer une borne au nombre de générateurs des
extensions d'anneaux et une borne au cardinal de la famille d'ouverts affines au-dessus
d'un ouvert affine donné de la base), où $\kappa$ est un cardinal, puis
construire un champ
$$
\Spacesstack_\kappa \longrightarrow (\Sch/S)_{fppf}
$$
exactement de la même manière que ci-dessus (à condition de veiller à ce que
$\Sch$ soit assez grande en fonction de $\kappa$).
\end{remark}





\section{Exemples de champs en groupoïdes}
\label{section-examples-stacks-in-groupoids}

\noindent
Les exemples précédents sont des champs qui ne sont pas des champs en
groupoïdes. Dans le reste de ce chapitre, nous donnons
des exemples issus de la géométrie algébrique de champs en groupoïdes.



\section{Le champ associé à un faisceau}
\label{section-stack-associated-to-sheaf}

\noindent
Soit $F : (\Sch/S)_{fppf}^{opp} \to \textit{Ensembles}$ un préfaisceau.
Nous obtenons une catégorie fibrée en ensembles
$$
p_F : \mathcal{S}_F \to (\Sch/S)_{fppf},
$$
voir
Catégories, exemple \ref{categories-example-presheaf}.
C'est un champ en ensembles si et seulement si $F$ est un faisceau, voir
Champs, lemme \ref{stacks-lemma-stack-in-setoids-characterize}.



\section{Le champ en groupoïdes des faisceaux quasi-cohérents de type fini}
\label{section-stack-in-groupoids-of-quasi-coherent-sheaves}

\noindent
Soit $p : \QCohstack_{fg} \to (\Sch/S)_{fppf}$ le champ
introduit à la
section \ref{section-stack-of-finitely-generated-quasi-coherent-sheaves}
(avec l'abus de notation qui y a été introduit).
Nous pouvons le transformer en un champ en groupoïdes
$p' : \QCohstack_{fg}' \to (\Sch/S)_{fppf}$ au moyen
du procédé de
Catégories, lemme \ref{categories-lemma-fibred-gives-fibred-groupoids},
voir
Champs, lemme \ref{stacks-lemma-stack-gives-stack-groupoids}.
Dans ce cas particulier, cela signifie simplement que $\QCohstack_{fg}'$ a
les mêmes objets que $\QCohstack_{fg}$, mais que ses morphismes sont les
couples $(f, g) : (U, \mathcal{F}) \to (U', \mathcal{F}')$
où $g$ est un isomorphisme $g : f^*\mathcal{F}' \to \mathcal{F}$.


\section{Le champ en groupoïdes des espaces algébriques de type fini}
\label{section-stack-in-groupoids-of-finite-type-spaces}

\noindent
Soit $p : \Spacesstack_{ft} \to (\Sch/S)_{fppf}$ le champ
introduit à la
section \ref{section-stack-of-finite-type-spaces}
(avec l'abus de notation qui y a été introduit).
Nous pouvons le transformer en un champ en groupoïdes
$p' : \Spacesstack_{ft}' \to (\Sch/S)_{fppf}$ au moyen
du procédé de
Catégories, lemme \ref{categories-lemma-fibred-gives-fibred-groupoids},
voir
Champs, lemme \ref{stacks-lemma-stack-gives-stack-groupoids}.
Dans ce cas particulier, cela signifie simplement que $\Spacesstack_{ft}'$
a les mêmes objets que $\Spacesstack_{ft}$, c'est-à-dire les morphismes de type fini
$X \to U$, où $X$ est un espace algébrique sur $S$ et $U$ un schéma
sur $S$. Mais les morphismes $(f, g) : X/U \to Y/V$ sont maintenant les
diagrammes commutatifs
$$
\xymatrix{
X \ar[d] \ar[r]_f & Y \ar[d] \\
U \ar[r]^g & V
}
$$
qui sont cartésiens.



\section{Champs quotients}
\label{section-quotient-stacks}

\noindent
Soit $(U, R, s, t, c)$ un groupoïde en espaces algébriques sur $S$.
Dans ce cas, le champ quotient
$$
[U/R] \longrightarrow (\Sch/S)_{fppf}
$$
est par construction un champ en groupoïdes, voir
Groupoïdes dans les espaces,
définition \ref{spaces-groupoids-definition-quotient-stack}.
Les faisceaux $\mathit{Isom}$ sont même
représentables par des espaces algébriques, voir
Amorçage, lemme \ref{bootstrap-lemma-quotient-stack-isom}.
Ces champs quotients sont d'une importance fondamentale dans la théorie des
champs algébriques.

\medskip\noindent
Un cas particulier de la construction précédente est le champ quotient
$$
[X/G] \longrightarrow (\Sch/S)_{fppf}
$$
associé à une donnée $(B, G/B, m, X/B, a)$. Ici,
\begin{enumerate}
\item $B$ est un espace algébrique sur $S$ ;
\item $(G, m)$ est un espace algébrique en groupes sur $B$ ;
\item $X$ est un espace algébrique sur $B$ ;
\item $a : G \times_B X \to X$ est une action de $G$ sur $X$ au-dessus de $B$.
\end{enumerate}
En effet, d'après
Groupoïdes dans les espaces,
définition \ref{spaces-groupoids-definition-quotient-stack},
le champ en groupoïdes $[X/G]$ est le
champ quotient $[X/G \times_B X]$ ci-dessus. Il convient
d'expliciter la catégorie $[X/G]$. Nous le ferons à la
section \ref{section-group-quotient-stacks}.



\section{Classification des torseurs}
\label{section-torsors}

\noindent
Nous allons expliquer avec soin plusieurs variantes de ce que peut
signifier l'étude du champ des torseurs sous un espace algébrique en groupes $G$
ou sous un faisceau de groupes $\mathcal{G}$.



\subsection{Torseurs sous un faisceau de groupes}
\label{subsection-torsors-sheaf}

\noindent
Soit $\mathcal{G}$ un faisceau de groupes sur $(\Sch/S)_{fppf}$.
Pour $U \in \Ob((\Sch/S)_{fppf})$, nous notons
$\mathcal{G}|_U$ la restriction de $\mathcal{G}$ à $(\Sch/U)_{fppf}$.
Nous définissons une catégorie $\mathcal{G}\textit{-Torsors}$ comme suit :
\begin{enumerate}
\item Un objet de $\mathcal{G}\textit{-Torsors}$ est un couple
$(U, \mathcal{F})$, où $U$ est un objet de $(\Sch/S)_{fppf}$
et $\mathcal{F}$ un torseur sous $\mathcal{G}|_U$, voir
Cohomologie sur les sites, définition \ref{sites-cohomology-definition-torsor}.
\item Un morphisme $(U, \mathcal{F}) \to (V, \mathcal{H})$ est donné
par un couple $(f, \alpha)$, où $f : U \to V$ est un morphisme de schémas
sur $S$, et $\alpha : f^{-1}\mathcal{H} \to \mathcal{F}$ un
isomorphisme de torseurs sous $\mathcal{G}|_U$.
\end{enumerate}
Ainsi, $\mathcal{G}\textit{-Torsors}$ est une catégorie et
$$
p : \mathcal{G}\textit{-Torsors} \longrightarrow (\Sch/S)_{fppf},
\quad
(U, \mathcal{F}) \longmapsto U
$$
est un foncteur. La catégorie fibre de $\mathcal{G}\textit{-Torsors}$
au-dessus de $U$ est la catégorie des torseurs sous $\mathcal{G}|_U$, qui est un groupoïde.

\begin{lemma}
\label{lemma-torsors-sheaf-stack-in-groupoids}
Moyennant un remplacement comme dans
Champs, remarque \ref{stacks-remark-stack-make-small},
le foncteur
$$
p : \mathcal{G}\textit{-Torsors} \longrightarrow (\Sch/S)_{fppf}
$$
définit un champ en groupoïdes sur $(\Sch/S)_{fppf}$.
\end{lemma}

\begin{proof}
La partie la plus difficile de la démonstration consiste à montrer
l'effectivité de la descente des objets.
Soit $\{U_i \to U\}_{i \in I}$ un recouvrement de $(\Sch/S)_{fppf}$.
Supposons donné, pour chaque $i$, un torseur $\mathcal{F}_i$
sous $\mathcal{G}|_{U_i}$ et, pour chaque $i, j \in I$, un isomorphisme
$\varphi_{ij} :
\mathcal{F}_i|_{U_i \times_U U_j} \to \mathcal{F}_j|_{U_i \times_U U_j}$
de torseurs sous $\mathcal{G}|_{U_i \times_U U_j}$,
satisfaisant une condition de cocycle convenable sur $U_i \times_U U_j \times_U U_k$.
Alors, d'après
Sites, section \ref{sites-section-glueing-sheaves},
nous obtenons sur $(\Sch/U)_{fppf}$ un faisceau $\mathcal{F}$
dont la restriction à chaque $U_i$ redonne $\mathcal{F}_i$ ainsi que
la donnée de descente. Par l'équivalence de catégories de
Sites, lemme \ref{sites-lemma-mapping-property-glue},
les morphismes d'action $\mathcal{G}|_{U_i} \times \mathcal{F}_i \to \mathcal{F}_i$
se recollent en un morphisme $a : \mathcal{G}|_U \times \mathcal{F} \to \mathcal{F}$.
Il reste à montrer que $a$ est une action et que $\mathcal{F}$ devient
un torseur sous $\mathcal{G}|_U$. Ces deux propriétés se vérifient localement
et résultent donc des propriétés correspondantes des actions
$\mathcal{G}|_{U_i} \times \mathcal{F}_i \to \mathcal{F}_i$.
Cela prouve l'effectivité de la descente des objets dans
$\mathcal{G}\textit{-Torsors}$.
Nous omettons certains détails.
\end{proof}



\subsection{Variante pour les torseurs sous un faisceau}
\label{subsection-variant-torsor-sheaf}

\noindent
On peut généraliser légèrement la construction de la
sous-section \ref{subsection-torsors-sheaf}.
Soient en effet $\mathcal{G} \to \mathcal{B}$ un morphisme de faisceaux
sur $(\Sch/S)_{fppf}$ et
$$
m :
\mathcal{G} \times_\mathcal{B} \mathcal{G}
\longrightarrow
\mathcal{G}
$$
une loi de groupe sur $\mathcal{G}/\mathcal{B}$. Autrement dit, le couple
$(\mathcal{G}, m)$ est un objet en groupes du topos
$\Sh((\Sch/S)_{fppf})/\mathcal{B}$. Voir
Sites, section \ref{sites-section-localize-topoi},
pour les localisations de topos.
Dans cette situation, nous pouvons définir une catégorie
$\mathcal{G}/\mathcal{B}\textit{-Torsors}$ comme suit
(nous utilisons le plongement de Yoneda pour considérer les schémas comme des faisceaux) :
\begin{enumerate}
\item Un objet de $\mathcal{G}/\mathcal{B}\textit{-Torsors}$ est un triplet
$(U, b, \mathcal{F})$, où
\begin{enumerate}
\item $U$ est un objet de $(\Sch/S)_{fppf}$ ;
\item $b : U \to \mathcal{B}$ est une section de $\mathcal{B}$ sur $U$ ;
\item $\mathcal{F}$ est un torseur sous $U \times_{b, \mathcal{B}}\mathcal{G}$
sur $U$.
\end{enumerate}
\item Un morphisme $(U, b, \mathcal{F}) \to (U', b', \mathcal{F}')$ est donné
par un couple $(f, g)$, où $f : U \to U'$ est un morphisme de schémas
sur $S$ tel que $b = b' \circ f$, et où
$g : f^{-1}\mathcal{F}' \to \mathcal{F}$ est un
isomorphisme de torseurs sous $U \times_{b, \mathcal{B}} \mathcal{G}$.
\end{enumerate}
Ainsi, $\mathcal{G}/\mathcal{B}\textit{-Torsors}$ est une catégorie et
$$
p :
\mathcal{G}/\mathcal{B}\textit{-Torsors}
\longrightarrow
(\Sch/S)_{fppf},
\quad
(U, b, \mathcal{F}) \longmapsto U
$$
est un foncteur. La catégorie fibre de
$\mathcal{G}/\mathcal{B}\textit{-Torsors}$
au-dessus de $U$ est la réunion disjointe, lorsque $b : U \to \mathcal{B}$ varie,
des catégories de torseurs sous $U \times_{b, \mathcal{B}} \mathcal{G}$ ;
c'est donc un groupoïde.

\medskip\noindent
Dans le cas particulier $\mathcal{B} = S$, on retrouve la catégorie
$\mathcal{G}\textit{-Torsors}$ introduite dans la
sous-section \ref{subsection-torsors-sheaf}.

\begin{lemma}
\label{lemma-variant-torsors-sheaf-stack-in-groupoids}
Moyennant un remplacement comme dans
Champs, remarque \ref{stacks-remark-stack-make-small},
le foncteur
$$
p :
\mathcal{G}/\mathcal{B}\textit{-Torsors}
\longrightarrow
(\Sch/S)_{fppf}
$$
définit un champ en groupoïdes sur $(\Sch/S)_{fppf}$.
\end{lemma}

\begin{proof}
Cette démonstration reprend celle du
lemme \ref{lemma-torsors-sheaf-stack-in-groupoids}.
Le lecteur est invité à lire d'abord cette dernière, dont
les notations sont moins lourdes.
La partie la plus difficile de la démonstration consiste à montrer
l'effectivité de la descente des objets. Soit $\{U_i \to U\}_{i \in I}$
un recouvrement de $(\Sch/S)_{fppf}$.
Supposons donné, pour chaque $i$, un couple $(b_i, \mathcal{F}_i)$
formé d'un morphisme $b_i : U_i \to \mathcal{B}$ et d'un
torseur $\mathcal{F}_i$ sous $U_i \times_{b_i, \mathcal{B}} \mathcal{G}$,
et supposons que, pour chaque $i, j \in I$,
on ait $b_i|_{U_i \times_U U_j} = b_j|_{U_i \times_U U_j}$ et
que soit donné un isomorphisme
$\varphi_{ij} :
\mathcal{F}_i|_{U_i \times_U U_j} \to \mathcal{F}_j|_{U_i \times_U U_j}$
de torseurs sous $(U_i \times_U U_j) \times_\mathcal{B} \mathcal{G}$,
satisfaisant une condition de cocycle convenable sur $U_i \times_U U_j \times_U U_k$.
Alors, d'après
Sites, section \ref{sites-section-glueing-sheaves},
nous obtenons sur $(\Sch/U)_{fppf}$ un faisceau $\mathcal{F}$
dont la restriction à chaque $U_i$ redonne $\mathcal{F}_i$ ainsi que
la donnée de descente. L'axiome de faisceau pour $\mathcal{B}$ montre
que les morphismes $b_i$ proviennent d'un unique morphisme $b : U \to \mathcal{B}$.
Par l'équivalence de catégories de
Sites, lemme \ref{sites-lemma-mapping-property-glue},
les morphismes d'action
$(U_i \times_{b_i, \mathcal{B}} \mathcal{G}) \times_{U_i} \mathcal{F}_i
\to \mathcal{F}_i$
se recollent en un morphisme
$(U \times_{b, \mathcal{B}} \mathcal{G}) \times \mathcal{F} \to \mathcal{F}$.
Il reste à montrer que c'est une action et que $\mathcal{F}$ devient
un torseur sous $U \times_{b, \mathcal{B}} \mathcal{G}$.
Ces deux propriétés se vérifient localement
et résultent donc des propriétés correspondantes des actions
sur les $\mathcal{F}_i$.
Cela prouve l'effectivité de la descente des objets dans
$\mathcal{G}/\mathcal{B}\textit{-Torsors}$.
Nous omettons certains détails.
\end{proof}



\subsection{Espaces principaux homogènes}
\label{subsection-principal-homogeneous-spaces}

\noindent
Soit $B$ un espace algébrique sur $S$.
Soit $G$ un espace algébrique en groupes sur $B$.
Nous définissons une catégorie $G\textit{-Principal}$ comme suit :
\begin{enumerate}
\item Un objet de $G\textit{-Principal}$ est un triplet $(U, b, X)$, où
\begin{enumerate}
\item $U$ est un objet de $(\Sch/S)_{fppf}$ ;
\item $b : U \to B$ est un morphisme sur $S$ ;
\item $X$ est un espace principal homogène sous $G_U$ sur $U$, où
$G_U = U \times_{b, B} G$.
\end{enumerate}
Voir
Groupoïdes dans les espaces,
définition \ref{spaces-groupoids-definition-principal-homogeneous-space}.
\item Un morphisme $(U, b, X) \to (U', b', X')$ est donné
par un couple $(f, g)$, où $f : U \to U'$ est un morphisme de schémas
sur $B$, et $g : X \to U \times_{f, U'} X'$ un
isomorphisme d'espaces principaux homogènes sous $G_U$.
\end{enumerate}
Ainsi, $G\textit{-Principal}$ est une catégorie et
$$
p : G\textit{-Principal} \longrightarrow (\Sch/S)_{fppf},
\quad
(U, b, X) \longmapsto U
$$
est un foncteur. La catégorie fibre de $G\textit{-Principal}$
au-dessus de $U$ est la réunion disjointe, lorsque $b : U \to B$ varie,
des catégories d'espaces principaux homogènes sous $U \times_{b, B} G$ ;
c'est donc un groupoïde.

\medskip\noindent
Dans le cas particulier $S = B$, les objets sont simplement les couples
$(U, X)$, où $U$ est un schéma sur $S$ et $X$ un espace principal homogène
sous $G_U$ sur $U$. De plus, les morphismes sont simplement les diagrammes
cartésiens
$$
\xymatrix{
X \ar[d] \ar[r]_g & X' \ar[d] \\
U \ar[r]^f & U'
}
$$
où $g$ est $G$-équivariant.

\begin{remark}
\label{remark-principal-stack-in-groupoids}
Nous conjecturons que, moyennant un remplacement comme dans
Champs, remarque \ref{stacks-remark-stack-make-small},
le foncteur
$$
p : G\textit{-Principal} \longrightarrow (\Sch/S)_{fppf}
$$
définit un champ en groupoïdes sur $(\Sch/S)_{fppf}$. Cela résulterait
de l'énoncé suivant : étant donnés
\begin{enumerate}
\item un recouvrement $\{U_i \to U\}_{i \in I}$ de $(\Sch/S)_{fppf}$ ;
\item un espace algébrique en groupes $H$ sur $U$ ;
\item pour chaque $i$, un espace principal homogène $X_i$ sous $H_{U_i}$
sur $U_i$ ;
\item des isomorphismes $H$-équivariants
$\varphi_{ij} : X_{i, U_i \times_U U_j} \to X_{j, U_i \times_U U_j}$
satisfaisant la condition de cocycle,
\end{enumerate}
il existe un espace principal homogène $X$ sous $H$ sur $U$
qui redonne $(X_i, \varphi_{ij})$. La méthode employée dans la démonstration de
Amorçage, lemme \ref{bootstrap-lemma-descent-torsor},
ramène ce problème à une question ensembliste ; le lecteur qui néglige
les questions ensemblistes « saura » donc que le résultat est vrai. On trouve dans
\url{https://math.columbia.edu/~dejong/wordpress/?p=591}
une suggestion pour aborder ce problème.
\end{remark}



\subsection{Variante pour les espaces principaux homogènes}
\label{subsection-variant-principal-homogeneous-spaces}

\noindent
Soit $S$ un schéma. Soit $B = S$.
Soit $G$ un schéma en groupes sur $B = S$.
Dans cette situation, nous pouvons définir une sous-catégorie pleine
$G\textit{-Principal-Schémas} \subset G\textit{-Principal}$
dont les objets sont les couples $(U, X)$, où $U$ est un objet de
$(\Sch/S)_{fppf}$ et $X \to U$ un espace principal homogène
sous $G$ sur $U$ qui est représentable, c'est-à-dire un schéma.

\medskip\noindent
En général, $G\textit{-Principal-Schémas}$ n'est pas
un champ en groupoïdes sur $(\Sch/S)_{fppf}$. Cela tient
à ce qu'il existe effectivement, en général, des espaces principaux homogènes
qui ne sont pas des schémas ; la descente des objets n'est donc pas effective
en général.



\subsection{Torseurs pour la topologie fppf}
\label{subsection-fppf-torsors}

\noindent
Soit $B$ un espace algébrique sur $S$.
Soit $G$ un espace algébrique en groupes sur $B$.
Nous définissons une catégorie $G\textit{-Torsors}$ comme suit :
\begin{enumerate}
\item Un objet de $G\textit{-Torsors}$ est un triplet $(U, b, X)$, où
\begin{enumerate}
\item $U$ est un objet de $(\Sch/S)_{fppf}$ ;
\item $b : U \to B$ est un morphisme ;
\item $X$ est un $G_U$-torseur fppf sur $U$, où $G_U = U \times_{b, B} G$.
\end{enumerate}
Voir
Groupoïdes dans les espaces,
définition \ref{spaces-groupoids-definition-principal-homogeneous-space}.
\item Un morphisme $(U, b, X) \to (U', b', X')$ est donné
par un couple $(f, g)$, où $f : U \to U'$ est un morphisme de schémas
sur $B$, et $g : X \to U \times_{f, U'} X'$ un
isomorphisme de $G_U$-torseurs.
\end{enumerate}
Ainsi, $G\textit{-Torsors}$ est une catégorie et
$$
p : G\textit{-Torsors} \longrightarrow (\Sch/S)_{fppf},
\quad
(U, a, X) \longmapsto U
$$
est un foncteur. La catégorie fibre de $G\textit{-Torsors}$
au-dessus de $U$ est la réunion disjointe, lorsque $b : U \to B$ varie,
des catégories de $U \times_{b, B} G$-torseurs fppf ;
c'est donc un groupoïde.

\medskip\noindent
Dans le cas particulier $S = B$, les objets sont simplement les couples
$(U, X)$, où $U$ est un schéma sur $S$ et $X$ un
$G_U$-torseur fppf sur $U$. De plus, les morphismes sont simplement les diagrammes
cartésiens
$$
\xymatrix{
X \ar[d] \ar[r]_g & X' \ar[d] \\
U \ar[r]^f & U'
}
$$
où $g$ est $G$-équivariant.

\begin{lemma}
\label{lemma-torsors-stack-in-groupoids}
Moyennant un remplacement comme dans
Champs, remarque \ref{stacks-remark-stack-make-small},
le foncteur
$$
p : G\textit{-Torsors} \longrightarrow (\Sch/S)_{fppf}
$$
définit un champ en groupoïdes sur $(\Sch/S)_{fppf}$.
\end{lemma}

\begin{proof}
La partie la plus difficile de la démonstration consiste à montrer l'effectivité
de la descente des objets ; c'est le contenu de
Amorçage, lemme \ref{bootstrap-lemma-descent-torsor}.
Nous omettons la démonstration des axiomes (1) et (2) de
Champs, définition \ref{stacks-definition-stack-in-groupoids}.
\end{proof}

\begin{lemma}
\label{lemma-compare-torsors}
Soit $B$ un espace algébrique sur $S$. Soit $G$ un espace algébrique en
groupes sur $B$. Notons $\mathcal{G}$, resp.\ $\mathcal{B}$, l'espace algébrique
$G$, resp.\ $B$, considéré comme un faisceau sur $(\Sch/S)_{fppf}$.
Le foncteur
$$
G\textit{-Torsors} \longrightarrow \mathcal{G}/\mathcal{B}\textit{-Torsors}
$$
qui associe au triplet $(U, b, X)$ le triplet
$(U, b, \mathcal{X})$, où $\mathcal{X}$ est $X$ considéré comme un faisceau,
est une équivalence de champs en groupoïdes sur $(\Sch/S)_{fppf}$.
\end{lemma}

\begin{proof}
Nous allons utiliser
Champs, lemme \ref{stacks-lemma-characterize-essentially-surjective-when-ff},
pour le prouver. Le foncteur est pleinement fidèle, puisque la catégorie des
espaces algébriques sur $S$ est une sous-catégorie pleine de la catégorie des
faisceaux sur $(\Sch/S)_{fppf}$.
En outre, tous les objets (des deux côtés) sont des torseurs localement triviaux,
de sorte que la condition (2) du lemme cité ci-dessus est satisfaite.
Le foncteur est donc une équivalence.
\end{proof}



\subsection{Variante pour les torseurs en topologie fppf}
\label{subsection-variant-fppf-torsors}

\noindent
Soit $S$ un schéma. Soit $B = S$.
Soit $G$ un schéma en groupes sur $B = S$.
Dans cette situation, nous pouvons définir une sous-catégorie pleine
$G\textit{-Torsors-Schémas} \subset G\textit{-Torsors}$
dont les objets sont les couples $(U, X)$, où $U$ est un objet de
$(\Sch/S)_{fppf}$ et $X \to U$ un
$G$-torseur fppf sur $U$ qui est représentable, c'est-à-dire un schéma.

\medskip\noindent
En général, $G\textit{-Torsors-Schémas}$ n'est pas
un champ en groupoïdes sur $(\Sch/S)_{fppf}$. Cela tient
à ce qu'il existe effectivement, en général, des $G$-torseurs fppf
qui ne sont pas des schémas ; la descente des objets n'est donc pas effective
en général.




\section{Quotients par des actions de groupes}
\label{section-group-quotient-stacks}

\noindent
Nous avons maintenant introduit assez de notations pour décrire
plus précisément les champs $[X/G]$ de la
section \ref{section-quotient-stacks}.

\begin{situation}
\label{situation-quotient-stack}
Dans cette situation,
\begin{enumerate}
\item $S$ est un schéma contenu dans $\Sch_{fppf}$ ;
\item $B$ est un espace algébrique sur $S$ ;
\item $(G, m)$ est un espace algébrique en groupes sur $B$ ;
\item $\pi : X \to B$ est un espace algébrique sur $B$ ;
\item $a : G \times_B X \to X$ est une action de $G$ sur $X$ au-dessus de $B$.
\end{enumerate}
\end{situation}

\noindent
Dans cette situation, nous construisons comme suit une catégorie $[[X/G]]$\footnote{La notation
$[[X/G]]$ avec doubles crochets sert à distinguer cette catégorie du
champ $[X/G]$ introduit précédemment. Dans la
proposition \ref{proposition-equal-quotient-stacks},
nous montrerons que les deux sont canoniquement équivalents.
Nous utiliserons ensuite la notation $[X/G]$ pour désigner indifféremment l'un ou l'autre.} :
\begin{enumerate}
\item Un objet de $[[X/G]]$ consiste en un quadruplet
$(U, b, P, \varphi : P \to X)$, où
\begin{enumerate}
\item $U$ est un objet de $(\Sch/S)_{fppf}$ ;
\item $b : U \to B$ est un morphisme sur $S$ ;
\item $P$ est un $G_U$-torseur fppf sur $U$, où $G_U = U \times_{b, B} G$ ;
\item $\varphi : P \to X$ est un morphisme $G$-équivariant qui s'insère
dans le diagramme commutatif
$$
\xymatrix{
P \ar[d] \ar[r]_{\varphi} & X \ar[d] \\
U \ar[r]^b & B
}
$$
\end{enumerate}
\item Un morphisme de $[[X/G]]$ est un couple
$(f, g) : (U, b, P, \varphi) \to (U', b', P', \varphi')$,
où $f : U \to U'$ est un morphisme de schémas sur $B$
et $g : P \to P'$ un morphisme $G$-équivariant au-dessus de $f$
qui induit un isomorphisme $P \cong U \times_{f, U'} P'$ et vérifie
$\varphi = \varphi' \circ g$.
Autrement dit, $(f, g)$ s'insère dans le diagramme commutatif
suivant :
$$
\xymatrix{
P \ar[d] \ar[rrrd]_\varphi \ar[r]^g & P' \ar[d] \ar[rrd]^{\varphi'} \\
U \ar[rrrd]_b \ar[r]^f & U' \ar[rrd]^{b'} & & X \ar[d] \\
& & & B
}
$$
\end{enumerate}
Ainsi, $[[X/G]]$ est une catégorie et
$$
p : [[X/G]] \longrightarrow (\Sch/S)_{fppf},
\quad
(U, b, P, \varphi) \longmapsto U
$$
est un foncteur. La catégorie fibre de $[[X/G]]$
au-dessus de $U$ est la réunion disjointe, lorsque $b \in \Mor_S(U, B)$ varie,
des $U \times_{b, B} G$-torseurs fppf $P$ munis d'un morphisme
$G$-équivariant vers $X$. Les catégories fibres de $[[X/G]]$ sont donc des groupoïdes.

\medskip\noindent
Remarquons que le foncteur
$$
[[X/G]] \longrightarrow G\textit{-Torsors},
\quad
(U, b, P, \varphi) \longmapsto (U, b, P)
$$
est un $1$-morphisme de catégories sur $(\Sch/S)_{fppf}$.

\begin{lemma}
\label{lemma-group-quotient-stack-in-groupoids}
Moyennant un remplacement comme dans
Champs, remarque \ref{stacks-remark-stack-make-small},
le foncteur
$$
p : [[X/G]] \longrightarrow (\Sch/S)_{fppf}
$$
définit un champ en groupoïdes sur $(\Sch/S)_{fppf}$.
\end{lemma}

\begin{proof}
La partie la plus difficile de la démonstration consiste à montrer l'effectivité
de la descente des objets. Supposons que $\{U_i \to U\}_{i \in I}$ soit un recouvrement de
$(\Sch/S)_{fppf}$. Soient
$\xi_i = (U_i, b_i, P_i, \varphi_i)$ des objets de $[[X/G]]$ au-dessus de $U_i$,
et soit $\varphi_{ij} : \text{pr}_0^*\xi_i \to \text{pr}_1^*\xi_j$
une donnée de descente. En appliquant le foncteur
$[[X/G]] \to G\textit{-Torsors}$ ci-dessus, nous obtenons en particulier une donnée de descente
sur les triplets $(U_i, b_i, P_i)$ pour le champ en groupoïdes
$G\textit{-Torsors}$. Nous avons vu que
$G\textit{-Torsors}$ est un champ en groupoïdes
(lemme \ref{lemma-torsors-stack-in-groupoids}).
Nous pouvons donc supposer que $b_i = b|_{U_i}$ pour un morphisme $b : U \to B$, et
que $P_i = U_i \times_U P$ pour un $G_U = U \times_{b, B} G$-torseur fppf
$P$ sur $U$. Les morphismes $\varphi_i$ sont compatibles
avec la donnée de descente canonique sur les restrictions $U_i \times_U P$
et définissent donc un morphisme $\varphi : P \to X$. (On peut par exemple
utiliser
Sites, lemme \ref{sites-lemma-mapping-property-glue},
ou
Descente sur les espaces,
lemme \ref{spaces-descent-lemma-fpqc-universal-effective-epimorphisms},
pour obtenir $\varphi$.)
Cela prouve l'effectivité de la descente des objets.
Nous omettons la démonstration des axiomes (1) et (2) de
Champs, définition \ref{stacks-definition-stack-in-groupoids}.
\end{proof}

\begin{proposition}
\label{proposition-equal-quotient-stacks}
Dans la
situation \ref{situation-quotient-stack},
il existe une équivalence canonique
$$
[X/G] \longrightarrow [[X/G]]
$$
de champs en groupoïdes sur $(\Sch/S)_{fppf}$.
\end{proposition}

\begin{proof}
Nous donnons tous les détails afin de nous assurer que les définitions
s'accordent exactement comme il le faut.
Rappelons que $[X/G]$ est le champ quotient
associé au groupoïde en espaces algébriques
$(X, G \times_B X, s, t, c)$, voir
Groupoïdes dans les espaces,
définition \ref{spaces-groupoids-definition-quotient-stack}.
Cela signifie que $[X/G]$ est la champification de la
catégorie fibrée en groupoïdes $[X/_{\!p}G]$ associée au foncteur
$$
(\Sch/S)_{fppf} \longrightarrow \textit{Groupoïdes},
\quad
U \longmapsto (X(U), G(U) \times_{B(U)} X(U), s, t, c),
$$
où $s(g, x) = x$, $t(g, x) = a(g, x)$ et
$c((g, x), (g', x')) = (m(g, g'), x')$. D'après la construction de
Catégories, exemple \ref{categories-example-functor-groupoids},
un objet de $[X/_{\!p}G]$ est un couple $(U, x)$ avec $x \in X(U)$,
et un morphisme $(f, g) : (U, x) \to (U', x')$ de $[X/_{\!p}G]$
est donné par un morphisme de schémas $f : U \to U'$ et un élément
$g \in G(U)$ tel que $a(g, x) = x' \circ f$.
Nous pouvons donc définir un $1$-morphisme de champs en groupoïdes
$$
F_p : [X/_{\!p}G] \longrightarrow [[X/G]]
$$
par les règles suivantes : sur les objets, nous posons
$$
F_p(U, x) =
(U, \pi \circ x, G \times_{B, \pi \circ x} U, a \circ (\text{id}_G \times x)).
$$
Cette définition a un sens, car le diagramme
$$
\xymatrix{
G \times_{B, \pi \circ x} U \ar[d] \ar[r]_{\text{id}_G \times x} &
G \times_{B, \pi} X \ar[r]_-a &
X \ar[d]^\pi \\
U \ar[rr]^{\pi \circ x} & & B
}
$$
est commutatif, et les deux flèches horizontales sont $G$-équivariantes si l'on considère
les produits fibrés comme des $G$-torseurs triviaux sur $U$, resp.\ sur $X$.
Sur un morphisme $(f, g) : (U, x) \to (U', x')$, nous posons
$F_p(f, g) = (f, R_{g^{-1}})$, où $R_{g^{-1}}$ désigne la translation à droite
par l'inverse de $g$. Plus précisément, le morphisme
$F_p(f, g) : F_p(U, x) \to F_p(U', x')$ est donné par le diagramme cartésien
$$
\xymatrix{
G \times_{B, \pi \circ x} U \ar[d] \ar[r]_{R_{g^{-1}}} &
G \times_{B, \pi \circ x'} U' \ar[d] \\
U \ar[r]^f & U'
}
$$
où, sur les points à valeurs dans $T$, $R_{g^{-1}}$ est défini par
$$
R_{g^{-1}}(g', u) = (m(g', i(g(u))), f(u)).
$$
Pour vérifier que cela convient, il faut établir l'égalité
$$
a \circ (\text{id}_G \times x)
=
a \circ (\text{id}_G \times x') \circ R_{g^{-1}},
$$
qui est vraie puisque le membre de droite, appliqué au point à valeurs dans $T$
$(g', u)$, donne bien l'égalité souhaitée
\begin{align*}
a((\text{id}_G \times x')(m(g', i(g(u))), f(u)))
& =
a(m(g', i(g(u))), x'(f(u))) \\
& =
a(g', a(i(g(u)), x'(f(u)))) \\
& =
a(g', x(u)),
\end{align*}
car $a(g, x) = x' \circ f$ et donc $a(i(g), x' \circ f) = x$.

\medskip\noindent
La propriété universelle de la champification donnée dans
Champs, lemme \ref{stacks-lemma-stackify-groupoids-universal-property},
fournit une extension canonique $F : [X/G] \to [[X/G]]$ du $1$-morphisme
$F_p$ ci-dessus. Montrons d'abord que $F$ est pleinement fidèle.
Comme la source et le but sont des champs en groupoïdes,
il suffit pour cela de montrer que $F$ identifie les faisceaux $\mathit{Isom}$.
Choisissons un schéma $U$ et des objets $\xi, \xi'$ de
$[X/G]$ au-dessus de $U$. Nous voulons montrer que
$$
F :
\mathit{Isom}_{[X/G]}(\xi, \xi')
\longrightarrow
\mathit{Isom}_{[[X/G]]}(F(\xi), F(\xi'))
$$
est un isomorphisme de faisceaux. Il suffit de travailler localement
sur $U$ ; nous pouvons donc supposer que $\xi, \xi'$ proviennent d'objets
$(U, x)$, $(U, x')$ de $[X/_{\!p}G]$ au-dessus de $U$. Cela résulte directement
de la construction de la champification et est aussi détaillé dans
Groupoïdes dans les espaces,
section \ref{spaces-groupoids-section-explicit-quotient-stacks}.
Soit en utilisant directement la description ci-dessus des morphismes de
$[X/_{\!p}G]$, soit en appliquant
Groupoïdes dans les espaces,
lemme \ref{spaces-groupoids-lemma-quotient-stack-morphisms},
nous voyons que, dans ce cas,
$$
\mathit{Isom}_{[X/G]}(\xi, \xi') =
U \times_{(x, x'), X \times_S X, (s, t)} (G \times_B X).
$$
Un point à valeurs dans $T$ de ce produit fibré correspond à un couple
$(u, g)$, où $u \in U(T)$ et $g \in G(T)$, tel que
$a(g, x \circ u) = x' \circ u$. (Cela implique en particulier
$\pi \circ x \circ u = \pi \circ x' \circ u$.)
D'autre part, un point à valeurs dans $T$
de $\mathit{Isom}_{[[X/G]]}(F(\xi), F(\xi'))$ correspond par définition
à un morphisme $u : T \to U$ tel que
$\pi \circ x \circ u = \pi \circ x' \circ u : T \to B$, ainsi qu'à un isomorphisme
$$
R :
G \times_{B, \pi \circ x \circ u} T
\longrightarrow
G \times_{B, \pi \circ x' \circ u} T
$$
de $G_T$-torseurs triviaux compatible avec les morphismes donnés vers $X$.
Puisque les torseurs sont triviaux, on a $R = R_{g^{-1}}$
(multiplication à droite) pour un certain $g \in G(T)$. La compatibilité avec les morphismes
$a \circ (1_G, x \circ u), a \circ (1_G, x' \circ u) : G \times_B T \to X$
équivaut à la condition $a(g, x \circ u) = x' \circ u$.
Nous obtenons donc l'égalité voulue des faisceaux $\mathit{Isom}$.

\medskip\noindent
Sachant maintenant que $F$ est pleinement fidèle, nous pouvons appliquer
Champs, lemme \ref{stacks-lemma-characterize-essentially-surjective-when-ff}.
Pour montrer que $F$ est une équivalence, il suffit donc
de montrer que les objets de $[[X/G]]$ appartiennent localement pour la topologie fppf à l'image essentielle
de $F$. Cela est clair, puisque les torseurs fppf sont localement triviaux pour la topologie fppf,
et la conclusion en résulte.
\end{proof}

\begin{lemma}
\label{lemma-classifying-stacks}
\begin{slogan}
Le champ classifiant d'un schéma en groupes ou d'un espace algébrique en groupes.
\end{slogan}
Soit $S$ un schéma. Soit $B$ un espace algébrique sur $S$.
Soit $G$ un espace algébrique en groupes sur $B$. Alors les champs
en groupoïdes
$$
[B/G],\quad
[[B/G]],\quad
G\textit{-Torsors},\quad
\mathcal{G}/\mathcal{B}\textit{-Torsors}
$$
sont tous canoniquement équivalents.
Si $G \to B$ est plat et localement
de présentation finie, ils sont aussi équivalents à
$G\textit{-Principal}$.
\end{lemma}

\begin{proof}
L'équivalence
$G\textit{-Torsors} \to \mathcal{G}/\mathcal{B}\textit{-Torsors}$
est donnée au lemme \ref{lemma-compare-torsors}.
L'équivalence $[B/G] \to [[B/G]]$ est donnée par la
proposition \ref{proposition-equal-quotient-stacks}.
En développant la définition de $[[B/G]]$ donnée à la
section \ref{section-group-quotient-stacks},
nous voyons que $[[B//G]] = G\textit{-Torsors}$.

\medskip\noindent
Supposons enfin $G \to B$ plat et localement de présentation finie.
Pour montrer que le foncteur naturel
$G\textit{-Torsors} \to G\textit{-Principal}$ est une équivalence,
il suffit de montrer que, pour un schéma $U$ sur $B$,
un espace principal homogène $X \to U$ sous $G_U$
est localement trivial pour la topologie fppf. Par notre définition des espaces principaux homogènes
(Groupoïdes dans les espaces,
définition \ref{spaces-groupoids-definition-principal-homogeneous-space}),
il existe un recouvrement fpqc $\{U_i \to U\}$ tel que
$U_i \times_U X \cong G \times_B U_i$ comme espaces algébriques sur $U_i$.
Il en résulte que $X \to U$ est surjectif, plat et localement de présentation
finie, voir
Descente sur les espaces, lemmes
\ref{spaces-descent-lemma-descending-property-surjective},
\ref{spaces-descent-lemma-descending-property-flat} et
\ref{spaces-descent-lemma-descending-property-locally-finite-presentation}.
Choisissons un schéma $W$ et un morphisme étale surjectif $W \to X$.
Il résulte alors de ce qui précède que $\{W \to U\}$ est un recouvrement fppf
tel que $X_W \to W$ admette une section. Ainsi, $X$ est un $G_U$-torseur fppf.
\end{proof}

\begin{remark}
\label{remark-X-mod-G-group}
Soit $S$ un schéma.
Soit $G$ un groupe abstrait.
Soit $X$ un espace algébrique sur $S$.
Soit $G \to \text{Aut}_S(X)$ un homomorphisme de groupes.
Dans cette situation, nous pouvons définir $[[X/G]]$ de façon analogue
à ce qui précède, comme suit :
\begin{enumerate}
\item Un objet de $[[X/G]]$ consiste en un triplet
$(U, P, \varphi : P \to X)$, où
\begin{enumerate}
\item $U$ est un objet de $(\Sch/S)_{fppf}$ ;
\item $P$ est un faisceau sur $(\Sch/U)_{fppf}$ muni
d'une action de $G$ qui en fait un torseur sous le faisceau constant
de valeur $G$ ;
\item $\varphi : P \to X$ est un morphisme $G$-équivariant de faisceaux.
\end{enumerate}
\item Un morphisme
$(f, g) : (U, P, \varphi) \to (U', P', \varphi')$
est donné par un morphisme de schémas $f : T \to T'$
et un isomorphisme $G$-équivariant
$g : P \to f^{-1}P'$ tel que $\varphi = \varphi' \circ g$.
\end{enumerate}
Exactement comme ci-dessus, nous obtenons un foncteur
$$
[[X/G]] \longrightarrow (\Sch/S)_{fppf}
$$
qui fait de $[[X/G]]$ un champ en groupoïdes sur $(\Sch/S)_{fppf}$.
Le faisceau constant $\underline{G}$ est représentable (pourvu que le cardinal de $G$ ne soit
pas trop grand) par $G_S$ sur $(\Sch/S)_{fppf}$,
et cette version de $[[X/G]]$ est équivalente au champ
$[[X/G_S]]$ introduit ci-dessus.
\end{remark}






\section{Le champ de Picard}
\label{section-picard-stack}

\noindent
Dans cette section, nous introduisons le champ de Picard en toute généralité.
Dans le chapitre sur Quot et Hilb, nous montrerons qu'il s'agit, sous des hypothèses
convenables, d'un champ algébrique, voir
Quot, section \ref{quot-section-picard-stack}.

\medskip\noindent
Soit $S$ un schéma.
Soit $\pi : X \to B$ un morphisme d'espaces algébriques sur $S$.
Nous définissons une catégorie $\Picardstack_{X/B}$ comme suit :
\begin{enumerate}
\item Un objet est un triplet $(U, b, \mathcal{L})$, où
\begin{enumerate}
\item $U$ est un objet de $(\Sch/S)_{fppf}$ ;
\item $b : U \to B$ est un morphisme sur $S$ ;
\item $\mathcal{L}$ est un faisceau inversible sur le changement de base
$X_U = U \times_{b, B} X$.
\end{enumerate}
\item Un morphisme $(f, g) : (U, b, \mathcal{L}) \to (U', b', \mathcal{L}')$
est donné par un morphisme de schémas $f : U \to U'$ sur $B$ et un
isomorphisme $g : f^*\mathcal{L}' \to \mathcal{L}$.
\end{enumerate}
Le composé de
$(f, g) : (U, b, \mathcal{L}) \to (U', b', \mathcal{L}')$
avec
$(f', g') : (U', b', \mathcal{L}') \to (U'', b'', \mathcal{L}'')$
est $(f' \circ f, g \circ f^*(g'))$.
Nous obtenons ainsi une catégorie $\Picardstack_{X/B}$ et
$$
p : \Picardstack_{X/B} \longrightarrow (\Sch/S)_{fppf},
\quad
(U, b, \mathcal{L}) \longmapsto U
$$
est un foncteur. La catégorie fibre de $\Picardstack_{X/B}$ au-dessus de $U$
est la réunion disjointe, lorsque $b \in \Mor_S(U, B)$ varie, des catégories
de faisceaux inversibles sur $X_U = U \times_{b, B} X$. Les catégories
fibres sont donc des groupoïdes.

\begin{lemma}
\label{lemma-picard-stack}
Moyennant un remplacement comme dans
Champs, remarque \ref{stacks-remark-stack-make-small},
le foncteur
$$
\Picardstack_{X/B} \longrightarrow (\Sch/S)_{fppf}
$$
définit un champ en groupoïdes sur $(\Sch/S)_{fppf}$.
\end{lemma}

\begin{proof}
Comme d'habitude, la partie la plus difficile consiste à montrer l'effectivité de la descente des objets.
Soit pour cela $\{U_i \to U\}$ un recouvrement de $(\Sch/S)_{fppf}$.
Soit $\xi_i = (U_i, b_i, \mathcal{L}_i)$ un objet de
$\Picardstack_{X/B}$ au-dessus de $U_i$, et soit
$\varphi_{ij} : \text{pr}_0^*\xi_i \to \text{pr}_1^*\xi_j$
une donnée de descente. Elle implique en particulier que les morphismes
$b_i$ sont les restrictions d'un morphisme $b : U \to B$.
Écrivons $X_U = U \times_{b, B} X$ et
$X_i = U_i \times_{b_i, B} X =
U_i \times_U U \times_{b, B} X = U_i \times_U X_U$.
Le faisceau $\mathcal{L}_i$ est un $\mathcal{O}_{X_i}$-module inversible.
La famille $\{X_i \to X_U\}$ est elle aussi un recouvrement fppf.
De plus, la donnée de descente $\varphi_{ij}$ fournit une
donnée de descente sur les faisceaux inversibles $\mathcal{L}_i$ relativement
au recouvrement fppf $\{X_i \to X_U\}$.
Par conséquent, d'après
Descente sur les espaces,
proposition \ref{spaces-descent-proposition-fpqc-descent-quasi-coherent},
nous obtenons un unique faisceau inversible $\mathcal{L}$ sur $X_U$
qui redonne $\mathcal{L}_i$ et les données de descente sur $X_i$.
Le triplet $(U, b, \mathcal{L})$ est donc l'objet de
$\Picardstack_{X/B}$ au-dessus de $U$ que nous cherchions.
Nous omettons les détails.
\end{proof}





\section{Exemples de champs d'inertie}
\label{section-examples-inertia}

\noindent
Voici quelques exemples de champs d'inertie.

\begin{example}
\label{example-inertia-stack-of-X-mod-G}
Soit $S$ un schéma. Soit $G$ un groupe commutatif.
Soit $X \to S$ un schéma sur $S$.
Soit $a : G \times X \to X$ une action de $G$ sur $X$.
Pour $g \in G$, nous notons $g : X \to X$ l'automorphisme correspondant.
Dans ce cas, le champ d'inertie de $[X/G]$ (voir
remarque \ref{remark-X-mod-G-group})
est donné par
$$
I_{[X/G]} = \coprod\nolimits_{g\in G} [X^g/G],
$$
où, pour un élément $g$ de $G$, le symbole $X^g$ désigne le
schéma $X^g = \{x \in X \mid g(x) = x\}$. Plus précisément,
$X^g$ est le produit
fibré
$$
X^g =  X \times_{(1, 1), X \times_S X, (g, 1)} X.
$$
En effet, pour tout $S$-schéma $T$, un
point à valeurs dans $T$ du champ d'inertie de $[X/G]$ consiste en un
triplet $(P/T, \phi, \alpha)$ formé d'un $G$-torseur fppf
$P \to T$, d'un morphisme $G$-équivariant
$\phi : P \to X$ et
d'un automorphisme $\alpha$ de $P$ sur $T$ tel que
$\phi \circ \alpha = \phi$.
Puisque $G$ est un faisceau de groupes \emph{commutatifs},
$\alpha$ est donné, localement pour la topologie fppf sur $T$,
par la multiplication par un élément $g$ de $G$.
La condition $\phi \circ \alpha = \phi$ signifie que $\phi$
se factorise par l'inclusion de $X^g$
dans $X$, c'est-à-dire que $\phi$ est le composé de cette inclusion avec un
morphisme $P \to X^g$.
La discussion précédente permet de définir un morphisme de catégories fibrées
$I_{[X/G]} \to \coprod_{g\in G} [X^g/G]$ sur les points à valeurs dans $T$ comme ci-dessus.
Nous omettons la vérification qu'il s'agit d'une équivalence.
\end{example}

\begin{example}
\label{example-inertia-stack-of-picard}
Soit $f : X \to S$ un morphisme de schémas.
Supposons que, pour tout $T \to S$, le changement de base $f_T : X_T \to T$
ait la propriété que le morphisme $\mathcal{O}_T \to f_{T, *}\mathcal{O}_{X_T}$
est un isomorphisme. (Cela implique que $f$ est
{\it cohomologiquement plat en dimension $0$} (insérer ici une référence ultérieure),
mais c'est une condition plus forte.) Considérons le champ de Picard $\Picardstack_{X/S}$, voir
section \ref{section-picard-stack}.
Les points de son champ d'inertie au-dessus d'un
$S$-schéma $T$ sont les couples $(\mathcal{L}, \alpha)$,
où $\mathcal{L}$ est un faisceau inversible
sur $X_T$ et $\alpha$ un automorphisme de ce faisceau inversible.
Autrement dit, nous pouvons considérer $\alpha$ comme un élément de
$H^0(X_T, \mathcal{O}_{X_T})^\times = H^0(T, \mathcal{O}_T^*)$
d'après notre hypothèse. Or $H^0(T, \mathcal{O}_T^*) = \mathbf{G}_{m, S}(T)$,
voir Groupoïdes, exemple \ref{groupoids-example-multiplicative-group}.
Le champ d'inertie de $\Picardstack_{X/S}$ est donc
$$
I_{\Picardstack_{X/S}} = \mathbf{G}_{m, S} \times_S \Picardstack_{X/S}
$$
en tant que champ sur $(\Sch/S)_{fppf}$.
\end{example}






\section{Champs de Hilbert finis}
\label{section-hilbert-d-stack}

\noindent
Nous formulons ce qui suit dans une généralité peut-être un peu plus grande
que strictement nécessaire. Fixons un $1$-morphisme
$$
F : \mathcal{X} \longrightarrow \mathcal{Y}
$$
de champs en groupoïdes sur $(\Sch/S)_{fppf}$. Pour chaque entier
$d \geq 1$, considérons une catégorie $\mathcal{H}_d(\mathcal{X}/\mathcal{Y})$
définie comme suit :
\begin{enumerate}
\item Un objet est un quintuplet $(U, Z, y, x, \alpha)$, où $U, Z$ sont des objets de
$(\Sch/S)_{fppf}$ et où $Z$ est fini et localement libre de degré
$d$ sur $U$, avec
$y \in \Ob(\mathcal{Y}_U)$, $x \in \Ob(\mathcal{X}_Z)$
et un isomorphisme $\alpha : y|_Z \to F(x)$\footnote{
Cela signifie que ces données fournissent, par le $2$-lemme de Yoneda
(Catégories, lemme \ref{categories-lemma-yoneda-2category}), un
diagramme $2$-commutatif
$$
\xymatrix{
(\Sch/Z)_{fppf} \ar[r]_-x \ar[d] & \mathcal{X} \ar[d]^F \\
(\Sch/U)_{fppf} \ar[r]^-y & \mathcal{Y}
}
$$
de champs en groupoïdes sur $(\Sch/S)_{fppf}$.
On peut aussi représenter $\alpha$ comme un $2$-morphisme
$$
\xymatrix{
(\Sch/Z)_{fppf}
\rrtwocell^{y \circ (Z \to U)}_{F \circ x}{\alpha} & &
\mathcal{Y}.
}
$$
}.
\item Un morphisme $(U, Z, y, x, \alpha) \to (U', Z', y', x', \alpha')$ est
donné par un morphisme de schémas $f : U \to U'$, un morphisme de schémas
$g : Z \to Z'$ qui induit un isomorphisme $Z \to Z' \times_U U'$,
et des isomorphismes $b : y \to f^*y'$, $a : x \to g^*x'$ induisant un diagramme
commutatif
$$
\xymatrix{
y|_Z \ar[rr]_\alpha \ar[d]_{b|_Z} & &
F(x) \ar[d]^{F(a)} \\
f^*y'|_Z \ar[rr]^{\alpha'} & &
F(g^*x') \\
}
$$
\end{enumerate}
Les définitions montrent immédiatement qu'il existe un foncteur d'oubli
canonique
$$
p :
\mathcal{H}_d(\mathcal{X}/\mathcal{Y})
\longrightarrow
(\Sch/S)_{fppf}
$$
qui associe au quintuplet $(U, Z, y, x, \alpha)$ le schéma $U$
et au morphisme
$(f, g, b, a) : (U, Z, y, x, \alpha) \to (U', Z', y', x', \alpha')$
le morphisme $f : U \to U'$.

\begin{lemma}
\label{lemma-hilbert-d-stack}
La catégorie $\mathcal{H}_d(\mathcal{X}/\mathcal{Y})$, munie
du foncteur $p$ ci-dessus, définit un champ en groupoïdes sur
$(\Sch/S)_{fppf}$.
\end{lemma}

\begin{proof}
Comme d'habitude, la partie la plus difficile consiste à montrer l'effectivité de la descente des objets.
Soit pour cela $\{U_i \to U\}$ un recouvrement de $(\Sch/S)_{fppf}$.
Soit $\xi_i = (U_i, Z_i, y_i, x_i, \alpha_i)$ un objet de
$\mathcal{H}_d(\mathcal{X}/\mathcal{Y})$ au-dessus de $U_i$, et soit
$\varphi_{ij} : \text{pr}_0^*\xi_i \to \text{pr}_1^*\xi_j$
une donnée de descente. Remarquons d'abord que $\varphi_{ij}$
induit une donnée de descente $(Z_i/U_i, \varphi_{ij})$, qui est effective d'après
Descente, lemme \ref{descent-lemma-affine}.
On obtient ainsi un schéma $Z/U$ fini et localement libre de degré $d$ d'après
Descente, lemme \ref{descent-lemma-descending-property-finite-locally-free}.
Désormais, nous identifions $Z_i$ à $Z \times_U U_i$.
Ensuite, les objets $y_i$ des catégories fibres $\mathcal{Y}_{U_i}$
descendent en un objet $y$ de $\mathcal{Y}_U$, puisque $\mathcal{Y}$ est un
champ en groupoïdes. De même, les objets $x_i$ des catégories fibres
$\mathcal{X}_{Z_i}$ descendent en un objet $x$ de $\mathcal{X}_Z$, puisque
$\mathcal{X}$ est un champ en groupoïdes. Enfin, les isomorphismes donnés
$$
\alpha_i :
(y|_Z)_{Z_i} = y_i|_{Z_i}
\longrightarrow
F(x_i) = F(x|_{Z_i})
$$
se recollent en un morphisme $\alpha : y|_Z \to F(x)$, car $\mathcal{Y}$
est un champ et $\mathit{Isom}_\mathcal{Y}(y|_Z, F(x))$ est donc
un faisceau. Nous omettons les détails.
\end{proof}

\begin{definition}
\label{definition-hilbert-d-stack}
Nous appellerons $\mathcal{H}_d(\mathcal{X}/\mathcal{Y})$
le {\it champ de Hilbert fini de degré $d$ de $\mathcal{X}$ sur $\mathcal{Y}$}
construit ci-dessus. Si $\mathcal{Y} = S$, nous écrivons
$\mathcal{H}_d(\mathcal{X}) = \mathcal{H}_d(\mathcal{X}/\mathcal{Y})$.
Si $\mathcal{X} = \mathcal{Y} = S$, nous le notons $\mathcal{H}_d$.
\end{definition}

\noindent
Étant donné $F : \mathcal{X} \to \mathcal{Y}$ comme ci-dessus, nous avons les
$1$-morphismes naturels suivants de champs en groupoïdes sur
$(\Sch/S)_{fppf}$ :
\begin{equation}
\label{equation-diagram-hilbert-d-stack}
\vcenter{
\xymatrix{
\mathcal{H}_d(\mathcal{X}) \ar[rd] &
\mathcal{H}_d(\mathcal{X}/\mathcal{Y}) \ar[d] \ar[l] \ar[r] &
\mathcal{Y} \\
& \mathcal{H}_d
}
}
\end{equation}
Chacune des flèches est donnée par un « foncteur d'oubli ».

\begin{lemma}
\label{lemma-faithful-hilbert}
Le $1$-morphisme
$\mathcal{H}_d(\mathcal{X}/\mathcal{Y}) \to \mathcal{H}_d(\mathcal{X})$
est fidèle.
\end{lemma}

\begin{proof}
Pour vérifier que
$\mathcal{H}_d(\mathcal{X}/\mathcal{Y}) \to \mathcal{H}_d(\mathcal{X})$
est fidèle, il suffit de le vérifier sur les catégories fibres.
Soient $\xi = (U, Z, y, x, \alpha)$ et $\xi' = (U, Z', y', x', \alpha')$
deux objets de $\mathcal{H}_d(\mathcal{X}/\mathcal{Y})$ au-dessus du
schéma $U$. Soient $(g, b, a), (g', b', a') : \xi \to \xi'$ deux morphismes
de la catégorie fibre de $\mathcal{H}_d(\mathcal{X}/\mathcal{Y})$ au-dessus de $U$.
Les images de ces morphismes dans $\mathcal{H}_d(\mathcal{X})$ coïncident
si et seulement si $g = g'$ et $a = a'$. Le diagramme commutatif
$$
\xymatrix{
y|_Z \ar[rr]_\alpha \ar[d]_{b|_Z, \ b'|_Z} & &
F(x) \ar[d]^{F(a) = F(a')} \\
y'|_Z \ar[rr]^-{\alpha'} & &
F(g^*x') = F((g')^*x') \\
}
$$
implique alors que $b|_Z = b'|_Z$. Puisque $Z \to U$ est fini et localement libre de degré
$d$, la famille $\{Z \to U\}$ est un recouvrement fppf ; ainsi $b = b'$.
\end{proof}












\begin{multicols}{2}[\section{Autres chapitres}]
\noindent
Préliminaires
\begin{enumerate}
\item \hyperref[introduction-section-phantom]{Introduction}
\item \hyperref[conventions-section-phantom]{Conventions}
\item \hyperref[sets-section-phantom]{Théorie des ensembles}
\item \hyperref[categories-section-phantom]{Catégories}
\item \hyperref[topology-section-phantom]{Topologie}
\item \hyperref[sheaves-section-phantom]{Faisceaux sur les espaces}
\item \hyperref[sites-section-phantom]{Sites et faisceaux}
\item \hyperref[stacks-section-phantom]{Champs}
\item \hyperref[fields-section-phantom]{Corps}
\item \hyperref[algebra-section-phantom]{Algèbre commutative}
\item \hyperref[brauer-section-phantom]{Groupes de Brauer}
\item \hyperref[homology-section-phantom]{Algèbre homologique}
\item \hyperref[derived-section-phantom]{Catégories dérivées}
\item \hyperref[simplicial-section-phantom]{Méthodes simpliciales}
\item \hyperref[more-algebra-section-phantom]{Compléments d'algèbre}
\item \hyperref[smoothing-section-phantom]{Lissage des morphismes d'anneaux}
\item \hyperref[modules-section-phantom]{Faisceaux de modules}
\item \hyperref[sites-modules-section-phantom]{Modules sur les sites}
\item \hyperref[injectives-section-phantom]{Objets injectifs}
\item \hyperref[cohomology-section-phantom]{Cohomologie des faisceaux}
\item \hyperref[sites-cohomology-section-phantom]{Cohomologie sur les sites}
\item \hyperref[dga-section-phantom]{Algèbre différentielle graduée}
\item \hyperref[dpa-section-phantom]{Algèbre à puissances divisées}
\item \hyperref[sdga-section-phantom]{Faisceaux différentiels gradués}
\item \hyperref[hypercovering-section-phantom]{Hyperrecouvrements}
\end{enumerate}
Schémas
\begin{enumerate}
\setcounter{enumi}{25}
\item \hyperref[schemes-section-phantom]{Schémas}
\item \hyperref[constructions-section-phantom]{Constructions de schémas}
\item \hyperref[properties-section-phantom]{Propriétés des schémas}
\item \hyperref[morphisms-section-phantom]{Morphismes de schémas}
\item \hyperref[coherent-section-phantom]{Cohomologie des schémas}
\item \hyperref[divisors-section-phantom]{Diviseurs}
\item \hyperref[limits-section-phantom]{Limites de schémas}
\item \hyperref[varieties-section-phantom]{Variétés}
\item \hyperref[topologies-section-phantom]{Topologies sur les schémas}
\item \hyperref[descent-section-phantom]{Descente}
\item \hyperref[perfect-section-phantom]{Catégories dérivées de schémas}
\item \hyperref[more-morphisms-section-phantom]{Compléments sur les morphismes}
\item \hyperref[flat-section-phantom]{Compléments sur la platitude}
\item \hyperref[groupoids-section-phantom]{Schémas en groupoïdes}
\item \hyperref[more-groupoids-section-phantom]{Compléments sur les schémas en groupoïdes}
\item \hyperref[etale-section-phantom]{Morphismes étales de schémas}
\end{enumerate}
Thèmes de théorie des schémas
\begin{enumerate}
\setcounter{enumi}{41}
\item \hyperref[chow-section-phantom]{Homologie de Chow}
\item \hyperref[intersection-section-phantom]{Théorie de l'intersection}
\item \hyperref[pic-section-phantom]{Schémas de Picard des courbes}
\item \hyperref[weil-section-phantom]{Théories cohomologiques de Weil}
\item \hyperref[adequate-section-phantom]{Modules adéquats}
\item \hyperref[dualizing-section-phantom]{Complexes dualisants}
\item \hyperref[duality-section-phantom]{Dualité pour les schémas}
\item \hyperref[discriminant-section-phantom]{Discriminants et différentes}
\item \hyperref[derham-section-phantom]{Cohomologie de de Rham}
\item \hyperref[local-cohomology-section-phantom]{Cohomologie locale}
\item \hyperref[algebraization-section-phantom]{Géométrie algébrique et formelle}
\item \hyperref[curves-section-phantom]{Courbes algébriques}
\item \hyperref[resolve-section-phantom]{Résolution des surfaces}
\item \hyperref[models-section-phantom]{Réduction semi-stable}
\item \hyperref[functors-section-phantom]{Foncteurs et morphismes}
\item \hyperref[equiv-section-phantom]{Catégories dérivées de variétés}
\item \hyperref[pione-section-phantom]{Groupes fondamentaux de schémas}
\item \hyperref[etale-cohomology-section-phantom]{Cohomologie étale}
\item \hyperref[crystalline-section-phantom]{Cohomologie cristalline}
\item \hyperref[proetale-section-phantom]{Cohomologie pro-étale}
\item \hyperref[relative-cycles-section-phantom]{Cycles relatifs}
\item \hyperref[more-etale-section-phantom]{Compléments de cohomologie étale}
\item \hyperref[trace-section-phantom]{La formule des traces}
\end{enumerate}
Espaces algébriques
\begin{enumerate}
\setcounter{enumi}{64}
\item \hyperref[spaces-section-phantom]{Espaces algébriques}
\item \hyperref[spaces-properties-section-phantom]{Propriétés des espaces algébriques}
\item \hyperref[spaces-morphisms-section-phantom]{Morphismes d'espaces algébriques}
\item \hyperref[decent-spaces-section-phantom]{Espaces algébriques décents}
\item \hyperref[spaces-cohomology-section-phantom]{Cohomologie des espaces algébriques}
\item \hyperref[spaces-limits-section-phantom]{Limites d'espaces algébriques}
\item \hyperref[spaces-divisors-section-phantom]{Diviseurs sur les espaces algébriques}
\item \hyperref[spaces-over-fields-section-phantom]{Espaces algébriques sur des corps}
\item \hyperref[spaces-topologies-section-phantom]{Topologies sur les espaces algébriques}
\item \hyperref[spaces-descent-section-phantom]{Descente et espaces algébriques}
\item \hyperref[spaces-perfect-section-phantom]{Catégories dérivées d'espaces}
\item \hyperref[spaces-more-morphisms-section-phantom]{Compléments sur les morphismes d'espaces}
\item \hyperref[spaces-flat-section-phantom]{Platitude sur les espaces algébriques}
\item \hyperref[spaces-groupoids-section-phantom]{Groupoïdes dans les espaces algébriques}
\item \hyperref[spaces-more-groupoids-section-phantom]{Compléments sur les groupoïdes dans les espaces}
\item \hyperref[bootstrap-section-phantom]{Amorçage}
\item \hyperref[spaces-pushouts-section-phantom]{Sommes amalgamées d'espaces algébriques}
\end{enumerate}
Thèmes de géométrie
\begin{enumerate}
\setcounter{enumi}{81}
\item \hyperref[spaces-chow-section-phantom]{Groupes de Chow des espaces}
\item \hyperref[groupoids-quotients-section-phantom]{Quotients de groupoïdes}
\item \hyperref[spaces-more-cohomology-section-phantom]{Compléments sur la cohomologie des espaces}
\item \hyperref[spaces-simplicial-section-phantom]{Espaces simpliciaux}
\item \hyperref[spaces-duality-section-phantom]{Dualité pour les espaces}
\item \hyperref[formal-spaces-section-phantom]{Espaces algébriques formels}
\item \hyperref[restricted-section-phantom]{Algébrisation des espaces formels}
\item \hyperref[spaces-resolve-section-phantom]{Résolution des surfaces revisitée}
\end{enumerate}
Théorie des déformations
\begin{enumerate}
\setcounter{enumi}{89}
\item \hyperref[formal-defos-section-phantom]{Théorie formelle des déformations}
\item \hyperref[defos-section-phantom]{Théorie des déformations}
\item \hyperref[cotangent-section-phantom]{Le complexe cotangent}
\item \hyperref[examples-defos-section-phantom]{Problèmes de déformation}
\end{enumerate}
Champs algébriques
\begin{enumerate}
\setcounter{enumi}{93}
\item \hyperref[algebraic-section-phantom]{Champs algébriques}
\item \hyperref[examples-stacks-section-phantom]{Exemples de champs}
\item \hyperref[stacks-sheaves-section-phantom]{Faisceaux sur les champs algébriques}
\item \hyperref[criteria-section-phantom]{Critères de représentabilité}
\item \hyperref[artin-section-phantom]{Axiomes d'Artin}
\item \hyperref[quot-section-phantom]{Espaces de Quot et de Hilbert}
\item \hyperref[stacks-properties-section-phantom]{Propriétés des champs algébriques}
\item \hyperref[stacks-morphisms-section-phantom]{Morphismes de champs algébriques}
\item \hyperref[stacks-limits-section-phantom]{Limites de champs algébriques}
\item \hyperref[stacks-cohomology-section-phantom]{Cohomologie des champs algébriques}
\item \hyperref[stacks-perfect-section-phantom]{Catégories dérivées de champs}
\item \hyperref[stacks-introduction-section-phantom]{Introduction aux champs algébriques}
\item \hyperref[stacks-more-morphisms-section-phantom]{Compléments sur les morphismes de champs}
\item \hyperref[stacks-geometry-section-phantom]{La géométrie des champs}
\end{enumerate}
Thèmes de théorie des espaces de modules
\begin{enumerate}
\setcounter{enumi}{107}
\item \hyperref[moduli-section-phantom]{Champs de modules}
\item \hyperref[moduli-curves-section-phantom]{Espaces de modules de courbes}
\end{enumerate}
Divers
\begin{enumerate}
\setcounter{enumi}{109}
\item \hyperref[examples-section-phantom]{Exemples}
\item \hyperref[exercises-section-phantom]{Exercices}
\item \hyperref[guide-section-phantom]{Guide bibliographique}
\item \hyperref[desirables-section-phantom]{Desiderata}
\item \hyperref[coding-section-phantom]{Style de programmation}
\item \hyperref[obsolete-section-phantom]{Obsolète}
\item \hyperref[fdl-section-phantom]{Licence de documentation libre GNU}
\item \hyperref[index-section-phantom]{Index généré automatiquement}
\end{enumerate}
\end{multicols}


\bibliography{my}
\bibliographystyle{amsalpha}

\end{document}
